% ============================================================
% I. GIỚI THIỆU  (~0,7 trang) — SƯỜN, viết sau Phần A/B
% Hai bài được chọn độc lập, KHÔNG cần ép liên kết.
% ============================================================
\section{Giới thiệu}\label{sec:intro}

\todo{Viết đoạn văn hoàn chỉnh theo các ý dưới.}

% Ý 1 — Bối cảnh chung ngắn gọn: đồ thị tri thức (KG) là cách biểu diễn tri thức có cấu trúc;
%       LLM/LVLM hiện nay vừa có thể dùng KG, vừa có thể giúp xây KG.
% Ý 2 — Bài thứ nhất (seminar): GraphVis [deng2024graphvis] — bài toán QA thường thức có KG hỗ trợ;
%       ý tưởng vẽ đồ thị con thành ảnh cho LVLM; huấn luyện hai giai đoạn. -> Phần A.
% Ý 3 — Bài thứ hai: SciMKG [lu2026scimkg] — dựng tự động KG giáo dục K12 với văn bản, ảnh,
%       video, âm thanh bằng chuỗi nhiều LLM (EVIA) và căn chỉnh đa phương thức. -> Phần B.
% Ý 4 — Nội dung báo cáo (gạch đầu dòng):
%       - trình bày chi tiết cơ chế hoạt động của từng hệ thống (dữ liệu vào/ra, mô hình dùng ở mỗi bước,
%         cách huấn luyện hoặc cách gọi mô hình), đối chiếu bài báo với mã nguồn công bố;
%       - phân tích ảnh hưởng của kích thước/độ phân giải ảnh đồ thị lên GraphVis;
%       - tái lập thực nghiệm của cả hai bài, log W&B công khai;
%       - nhận xét về tính nhất quán số liệu và khả năng tái lập.
% Ý 5 — Cấu trúc báo cáo: II nền tảng; III–VIII GraphVis; IX–XIII SciMKG; XIV kết luận.
